\setboolean{@twoside}{true} 
\clearpage
\pagestyle{dictionarystyle}

\chapter{Dictionary}
\label{chap:dictionary}

\subsection{A}

\begin{multicols*}{2}

    \dictentry{ABA}{CAS failure mode where a location returns to an old value 
    after intervening updates, so compare-and-swap succeeds incorrectly; 
    mitigated with tags or safe reclamation (Chapter~\ref{chap:memorymodel}).}

    \dictentry{ADL}{Argument Dependent Lookup}
        
    \dictentry{ADT}{Abstract Data Type}

    \dictentry{AF\_XDP}{Address Family XDP: Linux socket type for high-performance 
    packet I/O with zero-copy options; a kernel-bypass-adjacent path 
    (Chapter~\ref{chap:kernelbypass}).}

    \dictentry{Affinity}{binding a thread (or IRQ) to a CPU set so it does not 
    migrate; pairs with \texttt{isolcpus} and NUMA placement 
    (Chapter~\ref{chap:linuxprocesses}).}

    \dictentry{Allocate}{to assign, allot, distribute, or set apart for 
    a particular purpose}

    \dictentry{Allocator (C++)}{type that satisfies the Allocator requirements and 
    supplies \texttt{allocate}/\texttt{deallocate} for containers; PMR variants use 
    a runtime \texttt{memory\char`_resource} (Section~\ref{subsec:custom-allocators}).}

    \dictentry{Amortized}{cost of an operation averaged over a long sequence 
    (e.g.\ \texttt{vector::push\char`_back}); one step may be expensive, the 
    sequence is cheap (Section~\ref{subsec:amortized-vs-average}).}

    \dictentry{Arena}{preallocated memory region from which many objects are 
    carved; avoids hot-path \texttt{malloc} (Chapter~\ref{chap:memmag}).}

    \dictentry{Atomic}{Objects of atomic types are the only C++ objects 
    that are free from data races; that is, if one thread writes to an 
    atomic object while another thread reads from it, the behavior is 
    well-defined. In addition, accesses to atomic objects may establish 
    inter-thread synchronization and order non-atomic memory accesses 
    as specified by \textbf{\texttt{std::memory\char`_order}}.}
    
\end{multicols*}

\subsection{B}
\begin{multicols*}{2}

    \dictentry{BBO}{Best Bid and Offer: top-of-book prices (and often sizes); 
    the first level of the order book (Chapter~\ref{chap:algorithms}).}

    \dictentry{Big O}{asymptotic upper-bound notation for growth of time or 
    space with input size (Chapter~\ref{chap:algorithms}).}

    \dictentry{Black--Scholes}{classic continuous-time option pricing model; 
    basis for implied volatility (Chapter~\ref{chap:quantbasics}).}

    \dictentry{branch}{a mean for teams to develop features giving them 
    a separate workspace for their code. The branch to create are part 
    of the branching strategy, which could be different depending on 
    the specific svn used.}

    \dictentry{Busy polling}{spinning on I/O or a queue instead of sleeping 
    (e.g.\ \texttt{SO\char`_BUSY\char`_POLL}, epoll timeout 0); lower latency, 
    higher CPU use (Chapter~\ref{chap:lowlatency}).}

\end{multicols*}

\subsection{C}
\begin{multicols*}{2}

    \dictentry{Cache coherence}{protocol that keeps multiple CPU caches 
    consistent for shared lines (e.g.\ MESI); basis of false sharing 
    (Chapter~\ref{chap:cacheperformance}).}

    \dictentry{Cache line}{granularity of cache transfers and coherence 
    (often 64\,B); pad hot fields to avoid false sharing 
    (Chapter~\ref{chap:cacheperformance}).}

    \dictentry{CAS}{Compare-And-Swap: atomic read-modify-write that updates a 
    location only if it still holds an expected value 
    (\texttt{compare\char`_exchange\char`_*} in C++); basis of many lock-free 
    structures (Chapter~\ref{chap:memorymodel}).}

    \dictentry{CFS}{Completely Fair Scheduler: default Linux policy 
    (\texttt{SCHED\char`_OTHER}); migrates work across CPUs unless isolated 
    (Chapter~\ref{chap:linuxprocesses}).}

    \dictentry{Context switch}{saving/restoring thread or process CPU state 
    when the scheduler runs another task; costs $\mu$s and adds jitter 
    (Chapter~\ref{chap:linuxprocesses}).}

    \dictentry{constexpr}{a C++ specifier which declares that it is 
    possible to evaluate the value of the function or variable 
    at compile time.}

    \dictentry{CPU affinity}{see Affinity.}

\end{multicols*}

\subsection{D}

\begin{multicols*}{2}

    \dictentry{data race}{undefined behavior in C++ when two threads access the 
    same memory location, at least one access is a write, and the accesses are 
    not ordered by the memory model (Chapter~\ref{chap:memorymodel}).}

    \dictentry{deadlock}{any situation in which no member of some group 
    of entities can proceed because each waits for another member, 
    including itself, to take action, such as sending a message or, 
    more commonly, releasing a lock}

    \dictentry{Dependency injection}{design technique where a component receives 
    its dependencies from the outside (constructor/setter/framework) instead of 
    creating them itself; improves testability and decoupling.}

    \dictentry{DPDK}{Data Plane Development Kit: user-space packet processing 
    with poll-mode drivers; a common kernel-bypass stack 
    (Chapter~\ref{chap:kernelbypass}).}

    \dictentry{DSU}{Disjoint Set Union (Union--Find): nearly $O(1)$ amortized 
    connectivity queries (Section~\ref{subsec:pat-uf}).}

    \dictentry{duck typing}{a concept related to \textit{dynamic typing}, 
    where the type or the class of an object is less important than the 
    methods it defines}
    
    \dictentry{double dispatch}{choosing a method from the runtime types of 
    \emph{two} objects (e.g.\ Visitor): a single virtual call is not enough; 
    typically one virtual on the element then a virtual on the visitor.}
    
    \dictentry{dynamic binding}{selecting which function body runs at run time 
    (typically via \texttt{virtual} functions); opposed to static / early binding 
    (Section~\ref{sec:cpp-polymorphism}).}

\end{multicols*}

\subsection{E}
\begin{multicols*}{2}

    \dictentry{encapsulation}{preventing unauthorized access to some piece of information or functionality}

    \dictentry{Epoch reclaim}{safe memory reclamation for lock-free structures 
    by retiring nodes only after all threads leave a shared epoch 
    (Chapter~\ref{chap:memorymodel}).}

    \dictentry{EPOLLET}{edge-triggered \texttt{epoll}: notify on state change; 
    must drain until \texttt{EAGAIN} (Chapter~\ref{chap:eventdrivenio}).}

    \dictentry{epoll}{Linux scalable I/O event notification 
    (\texttt{epoll\char`_create}/\texttt{ctl}/\texttt{wait}); preferred over 
    \texttt{select}/\texttt{poll} for many file descriptors 
    (Chapter~\ref{chap:eventdrivenio}).}

\end{multicols*}

\subsection{F}
\begin{multicols*}{2}

    \dictentry{False sharing}{performance pathology where threads update different 
    variables on the same cache line, causing coherence ping-pong 
    (Chapter~\ref{chap:cacheperformance}).}

    \dictentry{First-touch}{NUMA page placement: a page is usually bound to the 
    node of the CPU that first writes it; pin threads before warm-up 
    (Chapter~\ref{chap:numa}).}

    \dictentry{FIX}{Financial Information eXchange: tag-value messaging standard 
    for many order/broker links; not usually the colocated matching-engine hot path 
    (Section~\ref{sec:market-wire-protocols}).}

    \dictentry{FOK}{Fill-or-Kill: execute the full quantity immediately or cancel 
    (Chapter~\ref{chap:quantbasics}).}

\end{multicols*}

\subsection{G}
\begin{multicols*}{2}

    \dictentry{Greeks}{option risk sensitivities (e.g.\ delta, gamma, theta, vega) 
    (Chapter~\ref{chap:quantbasics}).}

    \dictentry{GTC}{Good-Till-Cancel: order remains until filled or cancelled 
    (Chapter~\ref{chap:quantbasics}).}

\end{multicols*}

\subsection{H}

\begin{multicols*}{2}

    \dictentry{Happens-before}{memory-model relation that defines when one 
    operation's side effects are visible to another; established by 
    mutexes, atomics, etc.\ (Chapter~\ref{chap:memorymodel}).}

    \dictentry{Hazard pointers}{safe reclamation scheme: readers publish 
    pointers they may dereference; reclaimers skip those nodes 
    (Chapter~\ref{chap:memorymodel}).}

    \dictentry{Heap}{a large pool of memory used to satisfy memory 
    request by allocating part.}

    \dictentry{Huge pages}{large page sizes (e.g.\ 2\,MiB) that reduce TLB 
    pressure for big working sets (Chapter~\ref{chap:cacheperformance}).}

\end{multicols*}

\subsection{I}

\begin{multicols*}{2}

    \dictentry{Idiom}{a group of code fragments sharing an equivalent 
    semantic role (meaning), which recurs frequently across software 
    projects often expressing a special feature in one or more 
    programming languages or libraries.}

    \dictentry{Implied vol}{market price of an option expressed as the 
    Black--Scholes volatility that matches it (Chapter~\ref{chap:quantbasics}).}

    \dictentry{Interrupt coalescing}{NIC/driver (or OS) batches several packet 
    arrivals into fewer hard IRQs to raise throughput and cut interrupt rate. 
    That saves CPU but adds latency and jitter --- usually the wrong trade for 
    tick-to-trade paths (Chapter~\ref{chap:lowlatency}).}

    \dictentry{IOC}{Immediate-or-Cancel: take what is available now, cancel 
    the rest (Chapter~\ref{chap:quantbasics}).}

    \dictentry{I/O mux}{waiting on many file descriptors for readiness 
    (\texttt{select}/\texttt{poll}/\texttt{epoll}) 
    (Chapter~\ref{chap:eventdrivenio}).}

    \dictentry{irqaffinity}{kernel boot parameter setting the default CPU mask 
    for IRQs; keep defaults on housekeeping cores so isolated CPUs stay quiet 
    (Section~\ref{subsec:irq-affinity}).}

    \dictentry{irqbalance}{userspace daemon that reshuffles IRQ-to-CPU affinity 
    at runtime; disable or constrain on latency boxes or it undoes manual 
    \texttt{smp\char`_affinity} (Section~\ref{subsec:irq-affinity}).}

    \dictentry{isolcpus}{kernel boot parameter that removes CPUs from the general 
    scheduler balance set (scheduling exclusion); used with pinning for dedicated 
    HFT cores (Section~\ref{subsec:isolcpus}).}

    \dictentry{ITCH}{family of binary multicast market-data protocols 
    (order-book incremental updates) (Section~\ref{sec:market-wire-protocols}).}

    \dictentry{Iter. invalid.}{when a container operation makes iterators or 
    references unusable (e.g.\ \texttt{vector} reallocation) 
    (Section~\ref{sec:iterator-invalidation}).}

\end{multicols*}

\subsection{J}
\begin{multicols*}{2}

    \dictentry{Jitter}{variability of latency (not just the mean); isolation 
    and allocation-free paths target lower jitter (Chapter~\ref{chap:lowlatency}).}

\end{multicols*}

\subsection{K}
\begin{multicols*}{2}

    \dictentry{Kernel bypass}{moving packet path out of the kernel networking 
    stack (DPDK, OpenOnload, AF\_XDP, \ldots) to cut syscall/copy overhead 
    (Chapter~\ref{chap:kernelbypass}).}

\end{multicols*}

\subsection{L}
\begin{multicols*}{2}

    \dictentry{Locality}{spatial: nearby addresses used together; temporal: 
    same addresses reused soon; drives cache-friendly layout 
    (Chapter~\ref{chap:cacheperformance}).}

    \dictentry{Lock convoy}{many threads queueing on one mutex so throughput 
    collapses under contention (Chapter~\ref{chap:memorymodel}).}

    \dictentry{Lock-free}{algorithm that guarantees system-wide progress even if 
    some threads are delayed; typically built with atomics/CAS, not wait-free 
    for every thread (Chapter~\ref{chap:memorymodel}).}

    \dictentry{L1 (MD)}{market-data top-of-book (BBO) feed depth 
    (Chapter~\ref{chap:quantbasics}).}

    \dictentry{L2 (MD)}{market-data aggregated depth by price level 
    (Chapter~\ref{chap:quantbasics}).}

    \dictentry{L3 (MD)}{market-data full order-by-order / order-id book 
    (Chapter~\ref{chap:quantbasics}).}

\end{multicols*}

\subsection{M}
\begin{multicols*}{2}

    \dictentry{Matching engine}{venue component that matches orders and 
    publishes trades / book updates (Chapter~\ref{chap:quantbasics}).}

    \dictentry{Memory order}{C++ atomic ordering: \texttt{relaxed}, 
    \texttt{acquire}, \texttt{release}, \texttt{acq\char`_rel}, 
    \texttt{seq\char`_cst}; controls visibility and reordering 
    (Chapter~\ref{chap:memorymodel}).}

    \dictentry{MESI}{Modified / Exclusive / Shared / Invalid cache-coherence 
    protocol states (Chapter~\ref{chap:cacheperformance}).}

    \dictentry{method chaining}{API style where setters / builders return 
    \texttt{*this} (or a proxy) so calls can be written 
    \texttt{obj.setA(1).setB(2)}; common in the Named Parameter idiom.}

    \dictentry{mmap}{POSIX call that maps files or anonymous pages into a process 
    address space; basis for memory-mapped I/O and some shared memory 
    (Chapter~\ref{chap:memorymapping}).}

    \dictentry{mock}{test double that stands in for a dependency and records / 
    asserts how it was used (calls, arguments); contrasts with a stub, which 
    mostly returns canned data.}

    \dictentry{most vexing parsing}{C++ grammar ambiguity where a declaration 
    looks like an object definition but is parsed as a function declaration 
    (e.g.\ \texttt{TimeKeeper t(Timer());}); fix with extra parentheses or braces.}

    \dictentry{MPMC}{Multi-Producer Multi-Consumer queue; heavier than SPSC 
    (e.g.\ Michael--Scott) (Chapter~\ref{chap:memorymodel}).}

\end{multicols*}

\subsection{N}
\begin{multicols*}{2}

    \dictentry{Nagle}{TCP algorithm that coalesces small sends; disable with 
    \texttt{TCP\char`_NODELAY} for latency paths (Chapter~\ref{chap:linuxnetworking}).}

    \dictentry{NAPI}{Linux networking softirq path that polls the NIC after 
    an interrupt; can still preempt user threads 
    (Chapter~\ref{chap:kernelbypass}).}

    \dictentry{nohz\_full}{kernel boot option enabling adaptive-ticks / tickless 
    behaviour on listed CPUs (fewer scheduler timer interrupts when a single 
    task runs); often paired with \texttt{isolcpus} / \texttt{rcu\char`_nocbs} 
    (Section~\ref{subsec:nohz-full}).}

    \dictentry{NUMA}{Non-Uniform Memory Access: multi-socket systems where local 
    DRAM is faster than remote; place threads, memory, and NICs carefully 
    (Chapter~\ref{chap:numa}).}

\end{multicols*}

\subsection{O}

\begin{multicols*}{2}

    \dictentry{Object pool}{reuse preconstructed objects instead of 
    allocating on the hot path (Chapter~\ref{chap:memmag}).}

    \dictentry{OpenOnload}{Solarflare/Xilinx kernel-bypass stack that accelerates 
    sockets in user space (Chapter~\ref{chap:kernelbypass}).}

    \dictentry{object}{region of storage with associated semantics}

    \dictentry{Order book}{local structure of resting bids and asks by price 
    (and often order id); updated from market data (Chapter~\ref{chap:algorithms}).}

    \dictentry{Overloading}{a compile time polymorphism achieved 
    specifying more than one function with the same name but different 
    signature in the same scope.}
    
    \dictentry{Overridding}{the act of redefining a method of a base 
    class in a derived class, using the same name, return type and 
    parameters.}

\end{multicols*}

\subsection{P}

\begin{multicols*}{2}

    \dictentry{Page fault}{CPU trap when a virtual page is not mapped or not 
    present; major faults hit disk/IO and blow hot-path latency 
    (Chapter~\ref{chap:linuxprocesses}).}

    \dictentry{Parameterized type}
              {another definition for \highlight{class templates}.}

    \dictentry{perf}{Linux profiling tool (\texttt{perf record}/\texttt{stat}) 
    for cycles, cache misses, false sharing hunts 
    (Chapter~\ref{chap:cacheperformance}).}

    \dictentry{Pinning}{see Affinity.}

    \dictentry{PMD}{Poll-Mode Driver: NIC driver that busy-polls rings 
    (DPDK-style) instead of interrupt-driven RX (Chapter~\ref{chap:kernelbypass}).}

    \dictentry{PMR}{Polymorphic Memory Resource (\texttt{std::pmr}): runtime-selected 
    \texttt{memory\char`_resource} behind containers 
    (Section~\ref{subsec:custom-allocators}).}

    \dictentry{poll}{POSIX I/O multiplexer; improves on \texttt{select} but still 
    scales worse than \texttt{epoll} for large FD sets 
    (Chapter~\ref{chap:eventdrivenio}).}

    \dictentry{Polymorphysm}
              {ability of different objects to be accessed by a common 
              interface.}
                      
\end{multicols*}

%~ \subsection{Q}
%~ \begin{multicols*}{2}\end{multicols*}

\subsection{R}

\begin{multicols*}{2}
    
    \dictentry{Race Condition}{undesirable situation which occurs when 
    two instructions from different threads access the same memory 
    location, at least one of these accesses is a write and there is no 
    synchronization that is mandating any particular order among these 
    accesses.}
    
    \dictentry{RAII}{acronym for Resource Acquisition Is Initialization.}
    
    \dictentry{RBV optimization}{see~\cite{isocpp-rbv-optimization}.}

    \dictentry{RCU}{Read-Copy-Update: Linux synchronization for read-heavy shared 
    data; writers defer freeing old versions via callbacks. Those callbacks can 
    jitter isolated CPUs unless offloaded with \texttt{rcu\char`_nocbs} 
    (Section~\ref{subsec:rcu-nocbs}).}

    \dictentry{rcu\_nocbs}{kernel boot option that offloads RCU callback processing 
    from listed CPUs onto housekeeping cores (Section~\ref{subsec:rcu-nocbs}).}

    \dictentry{Reactor}{event-loop style: demultiplex I/O readiness then 
    dispatch handlers (typical \texttt{epoll} server) 
    (Chapter~\ref{chap:eventdrivenio}).}

    \dictentry{RTTI}{acronym for Run-Time Type Information: language mechanism 
    that exposes dynamic type of polymorphic objects (e.g.\ \texttt{typeid}, 
    \texttt{dynamic\char`_cast}). Requires a polymorphic type (usually a virtual 
    function); typically implemented via type info reachable from the vtable. 
    Can be disabled with compiler flags (\texttt{-fno-rtti}) to save space/cost.}

    \dictentry{RVO}{acronym for Return Value Optimization.}

\end{multicols*}

\subsection{S}
\begin{multicols*}{2}

    \dictentry{SBE}{Simple Binary Encoding: compact, schema-driven binary 
    messages for market data / OE (Section~\ref{sec:market-wire-protocols}).}

    \dictentry{SCHED\_FIFO}{Linux real-time FIFO scheduling policy; use carefully 
    on dedicated cores (Chapter~\ref{chap:linuxprocesses}).}

    \dictentry{SCHED\_RR}{Linux real-time round-robin scheduling policy 
    (Chapter~\ref{chap:linuxprocesses}).}

    \dictentry{select}{classic POSIX I/O multiplexer with an FD-count limit and 
    $O(n)$ scans; prefer \texttt{poll}/\texttt{epoll} for scale 
    (Chapter~\ref{chap:eventdrivenio}).}

    \dictentry{Seq. consistency}{strongest C++ atomic default 
    (\texttt{memory\char`_order\char`_seq\char`_cst}): a single total order of 
    all seq\_cst ops (Chapter~\ref{chap:memorymodel}).}

    \dictentry{smp\_affinity}{per-IRQ CPU mask under 
    \texttt{/proc/irq/<n>/smp\char`_affinity} (or 
    \texttt{smp\char`_affinity\char`_list}); steer hard IRQs (e.g.\ NIC queues) 
    onto housekeeping cores (Section~\ref{subsec:irq-affinity}).}

    \dictentry{softirq}{deferred kernel interrupt work (e.g.\ networking) that 
    can run on a CPU and add jitter (Chapter~\ref{chap:kernelbypass}).}

    \dictentry{Spinlock}{busy-wait lock for very short critical sections; 
    burns CPU while waiting (Chapter~\ref{chap:memorymodel}).}

    \dictentry{Spread}{difference between best ask and best bid 
    (Chapter~\ref{chap:quantbasics}).}
    
    \dictentry{static typing}{types are checked at compile time; variables and 
    expressions have types known to the compiler (as in C++), opposed to 
    dynamic typing where much checking happens at run time.}
    
    \dictentry{SFINAE}{acronym for \textbf{Substitution} \textbf{Failure} 
    \textbf{Is} \textbf{Not} \textbf{An} \textbf{Error}, a technique 
    used in template programming which discards the specialization from 
    the overloaded set avoiding a compile error, when substituting the 
    explicitly specified or deduced type for the template parameter fails.}

    \dictentry{SPSC}{Single-Producer Single-Consumer queue (often a ring buffer): 
    lock-free handoff between exactly two threads 
    (Section~\ref{sec:spsc-ring-buffer}).}
    
    \dictentry{Stack}{LIFO-ordered contiguous region of temporary memory 
    whose size is known to the compiler, which manages to allocate and 
    deallocate blocks of memory for the stack variables like functions' 
    local data. . If a region of memory lies on the thread's stack, 
    that memory is said to have been allocated on the stack, i.e. 
    stack-based memory allocation (SBMA).}
    
    \dictentry{Stub}{test double that provides canned answers to calls; lighter 
    than a mock (which also verifies interaction). In broader software slang, 
    also an unfinished placeholder implementation.}

    \dictentry{Syscall}{trap from user space into the kernel (e.g.\ 
    \texttt{recv}, \texttt{epoll\char`_wait}); mode switch cost motivates 
    bypass (Chapter~\ref{chap:kernelbypass}).}

\end{multicols*}


\newpage{}
\subsection{T}
\begin{multicols*}{2}

    \dictentry{Tail latency}{high percentiles of latency (p99 / p999), not the 
    mean; HFT cares about tails and jitter (Chapter~\ref{chap:lowlatency}).}

    \dictentry{TCP\_NODELAY}{socket option that disables Nagle coalescing for 
    lower send latency (Chapter~\ref{chap:linuxnetworking}).}

    \dictentry{template function}
              {instantiation of a function template}

    \dictentry{Test-driven Development}
              {\href{https://en.wikipedia.org/wiki/Test-driven_development}{wikipedia}}

    \dictentry{tick-to-trade}{end-to-end latency from receiving a market event 
    to sending an order (Chapter~\ref{chap:lowlatency}).}

    \dictentry{TLB}{Translation Lookaside Buffer: CPU cache of 
    virtual$\rightarrow$physical address translations 
    (Chapter~\ref{chap:cacheperformance}).}

    \dictentry{Type erasure}{hiding a concrete type behind a uniform interface 
    (virtuals, \texttt{std::function}/\texttt{any}, or a custom eraser); related 
    but not identical to opaque \texttt{void*} payloads 
    (Section~\ref{subsec:type-erasure}).}

\end{multicols*}

\subsection{U}
\begin{multicols*}{2}

    \dictentry{Union--Find}{see DSU.}

\end{multicols*}

\subsection{V}
\begin{multicols*}{2}
    
    \dictentry{virtual member function}{a mean to allow derived classes 
    to replace the implementation provided by the base class. The 
    compiler makes sure to call the correct replacement when 
    the object is of the derived class and also when the object is 
    accessed through a pointer to the base rather than to the derived 
    class.}
    
\end{multicols*}

\subsection{W}
\begin{multicols*}{2}

    \dictentry{Wait-free}{stronger than lock-free: every thread completes 
    its operation in a bounded number of its own steps 
    (Chapter~\ref{chap:memorymodel}).}

\end{multicols*}

\subsection{X}
\begin{multicols*}{2}

    \dictentry{XDP}{eXpress Data Path: early in-kernel packet hook; 
    AF\_XDP exposes XDP rings to user space (Chapter~\ref{chap:kernelbypass}).}

\end{multicols*}

%~ \subsection{Y}
%~ \begin{multicols*}{2}\end{multicols*}

\subsection{Z}
\begin{multicols*}{2}

    \dictentry{Zero-copy}{I/O path that avoids copying packet/payload bytes 
    between kernel and user buffers (Chapter~\ref{chap:kernelbypass}).}

\end{multicols*}



\setboolean{@twoside}{false} 
\pagestyle{basicstyle}
